% ============================================================================
%  ESTADÍSTICA INFERENCIAL — Sesión 15
%  Semana 15: 16 al 21 de noviembre de 2026   (lunes 16: festivo)
%  3.5 Correlación y regresión lineal · 3.6 Validación de modelos
%  Cierre de la Unidad 3 y repaso para el Examen 3
% ============================================================================
\documentclass[aspectratio=169,10pt]{beamer}

\usepackage[utf8]{inputenc}
\IfFileExists{spanish.ldf}{\usepackage[spanish,es-noshorthands]{babel}}{}
\usepackage{amsmath,amssymb}
\usepackage{booktabs}
\usepackage{tikz}
\usepackage{listings}

\usetheme{Uniminuto}

\uniasignatura{Estadística inferencial}
\uniprofesor{Andrés Felipe Puerta Velez}
\unicorreo{andres.puerta-v@uniminuto.edu.co}
\uninrc{CBASBL061}
\unisemana{Sesión 15 — Semana 15 · 16 al 21 de noviembre de 2026}

\newcommand{\listatiny}{\setbeamerfont{itemize/enumerate body}{size=\tiny}\tiny}

\newcommand{\fx}[1]{{\ttfamily\color{uniAzul}#1}}

% ------------------------------------------------------------- CÓDIGO R ----
\lstdefinestyle{Rcode}{
  language=R,
  basicstyle=\ttfamily\scriptsize,
  keywordstyle=\color{uniAzul}\bfseries,
  commentstyle=\color{uniGris}\itshape,
  stringstyle=\color{uniTitulo},
  showstringspaces=false, numbers=none, breaklines=true,
  columns=fullflexible,
  backgroundcolor=\color{uniCremaClaro},
  frame=leftline, framerule=1.2pt, rulecolor=\color{uniNaranja},
  xleftmargin=0.35cm, framexleftmargin=0.35cm,
  framexrightmargin=0.15cm, framextopmargin=2pt, framexbottommargin=2pt,
  aboveskip=0.15cm, belowskip=0.15cm
}
\lstdefinestyle{Rmini}{
  language=R,
  basicstyle=\ttfamily\tiny,
  keywordstyle=\color{uniAzul}\bfseries,
  commentstyle=\color{uniGris}\itshape,
  stringstyle=\color{uniTitulo},
  showstringspaces=false, numbers=none, breaklines=true,
  columns=fullflexible,
  backgroundcolor=\color{uniCremaClaro},
  frame=leftline, framerule=1.2pt, rulecolor=\color{uniNaranja},
  xleftmargin=0.35cm, framexleftmargin=0.35cm,
  framexrightmargin=0.15cm, framextopmargin=2pt, framexbottommargin=2pt,
  aboveskip=0.12cm, belowskip=0.12cm
}
\lstdefinestyle{Rout}{
  basicstyle=\ttfamily\tiny\color{black},
  showstringspaces=false, numbers=none, breaklines=true,
  columns=fullflexible,
  backgroundcolor=\color{uniGris!10},
  frame=leftline, framerule=1.2pt, rulecolor=\color{uniGris},
  xleftmargin=0.35cm, framexleftmargin=0.35cm,
  framexrightmargin=0.15cm, framextopmargin=2pt, framexbottommargin=2pt,
  aboveskip=0.12cm, belowskip=0.12cm
}

\begin{document}

\portada

\begin{frame}{Orden del día}
  \begin{enumerate}
    \item Llamado a lista y revisión de los ejercicios propuestos
    \item Repaso: dónde quedamos
    \item \azul{Tema 3.5: correlación y regresión lineal}
    \item \azul{Tema 3.6: validación de modelos}
    \item Ejercicios en clase
    \item \destacar{Repaso para el Examen 3}
  \end{enumerate}
\end{frame}

\begin{frame}{Repaso: dónde quedamos}
  Todo el curso ha trabajado con \azul{una} variable: estimar su media, su
  proporción, contrastar hipótesis sobre ella.

  \vspace{0.25cm}
  \begin{alertblock}{El último paso}
    \footnotesize
    Hoy trabajamos con \destacar{dos} variables a la vez y respondemos dos
    preguntas nuevas: ¿están relacionadas y qué tan fuerte? (correlación) y
    ¿puedo usar una para predecir la otra? (regresión).
  \end{alertblock}

  \vspace{0.15cm}
  {\footnotesize\color{uniGris} Con esto se cierra la unidad 3 y el curso. El
  \destacar{Examen 3} es la próxima semana: 23 al 28 de noviembre.}
\end{frame}

\seccion{3.5 Correlación}

\begin{frame}{Coeficiente de correlación de Pearson}
  \begin{block}{Fórmula}
    \[ r=\frac{n\sum xy-\sum x\sum y}
       {\sqrt{\left[n\sum x^{2}-(\sum x)^{2}\right]
              \left[n\sum y^{2}-(\sum y)^{2}\right]}} \]
  \end{block}

  \vspace{0.2cm}
  \begin{columns}[T,onlytextwidth]
    \begin{column}{0.48\textwidth}
      \begin{exampleblock}{Qué mide}
        \footnotesize
        La fuerza y el sentido de la relación \azul{lineal}. Siempre
        $-1\le r\le 1$.
      \end{exampleblock}
    \end{column}
    \begin{column}{0.48\textwidth}
      \begin{alertblock}{Cómo se lee}
        \footnotesize
        $r>0$: crecen juntas. $r<0$: una sube y la otra baja. $r\approx 0$: no
        hay relación \emph{lineal} (podría haber otra).
      \end{alertblock}
    \end{column}
  \end{columns}

  \vspace{0.2cm}
  \footnotesize
  \centering
  $|r|<0{,}3$ débil \quad · \quad $0{,}3$ a $0{,}7$ moderada \quad · \quad
  $>0{,}7$ fuerte
\end{frame}

\begin{frame}{Correlación no es causalidad}
  \begin{alertblock}{La advertencia más importante de la sesión}
    \footnotesize
    Que dos variables se muevan juntas no significa que una cause la otra. Un
    $r$ alto puede deberse a:
    \begin{itemize}
      \item Una \azul{tercera variable} que afecta a ambas (el tamaño de la
            ciudad explica a la vez ventas y número de empleados).
      \item Coincidencia en muestras pequeñas.
      \item Causalidad \azul{inversa}: quizá $y$ causa $x$.
    \end{itemize}
  \end{alertblock}

  \vspace{0.2cm}
  \begin{block}{Para afirmar causalidad hace falta más}
    \footnotesize
    Un diseño experimental con asignación aleatoria, o al menos control de las
    variables que podrían confundir. La estadística descriptiva de la relación
    no basta.
  \end{block}
\end{frame}

\seccion{3.5 Regresión lineal simple}

\begin{frame}{La recta de regresión}
  \begin{block}{Modelo}
    \[ \hat{y}=b_0+b_1 x \qquad
       b_1=\frac{n\sum xy-\sum x\sum y}{n\sum x^{2}-(\sum x)^{2}} \qquad
       b_0=\bar{y}-b_1\bar{x} \]
  \end{block}

  \vspace{0.2cm}
  \begin{columns}[T,onlytextwidth]
    \begin{column}{0.48\textwidth}
      \begin{exampleblock}{Interpretación de $b_1$}
        \footnotesize
        Cuánto cambia $y$ en promedio por cada unidad que aumenta $x$. Es la
        pendiente.
      \end{exampleblock}
    \end{column}
    \begin{column}{0.48\textwidth}
      \begin{exampleblock}{Interpretación de $b_0$}
        \footnotesize
        Valor estimado de $y$ cuando $x=0$. A veces no tiene sentido práctico:
        depende de si $x=0$ es un valor posible.
      \end{exampleblock}
    \end{column}
  \end{columns}

  \vspace{0.2cm}
  {\footnotesize\color{uniGris} Los coeficientes se obtienen por
  \azul{mínimos cuadrados}: la recta que hace mínima la suma de los cuadrados
  de los residuos $e_i=y_i-\hat{y}_i$.}
\end{frame}

\begin{frame}{Ejercicio 1 — en clase}
  \begin{alertblock}{En parejas (12 minutos)}
    Una empresa registra la inversión en publicidad ($x$, millones) y las
    ventas ($y$, millones) de 8 meses:

    \vspace{0.2cm}
    \centering\footnotesize
    \begin{tabular}{@{}lcccccccc@{}}
      \toprule
      $x$ & 2 & 4 & 5 & 7 & 8 & 10 & 12 & 14 \\
      $y$ & 35 & 40 & 52 & 55 & 68 & 70 & 78 & 88 \\
      \bottomrule
    \end{tabular}

    \vspace{0.2cm}
    \raggedright
    Sumas: $\sum x=62$, $\sum y=486$, $\sum xy=4287$, $\sum x^{2}=598$,
    $\sum y^{2}=31906$.

    \vspace{0.15cm}
    \begin{enumerate}
      \item Calcule $r$ e interprételo.
      \item Obtenga la recta $\hat{y}=b_0+b_1x$.
      \item Prediga las ventas con una inversión de 9 millones.
    \end{enumerate}
  \end{alertblock}
\end{frame}

\begin{frame}{Ejercicio 1 — solución}
  \begin{columns}[T,onlytextwidth]
    \begin{column}{0.52\textwidth}
      \begin{block}{Cálculos}
        \footnotesize
        \[ b_1=\frac{8(4287)-62(486)}{8(598)-62^{2}}
              =\frac{34296-30132}{4784-3844}=\azul{4{,}4298} \]
        \[ b_0=\frac{486}{8}-4{,}4298\left(\frac{62}{8}\right)=\azul{26{,}4191} \]
        \[ r=\azul{0{,}9840} \qquad r^{2}=\azul{0{,}9682} \]
      \end{block}
    \end{column}
    \begin{column}{0.44\textwidth}
      \begin{exampleblock}{Recta y predicción}
        \footnotesize
        \[ \hat{y}=26{,}42+4{,}43\,x \]
        \[ \hat{y}(9)=26{,}42+4{,}43(9)=\azul{66{,}29} \]
      \end{exampleblock}
      \begin{alertblock}{Interpretación}
        \footnotesize
        Correlación positiva \azul{fuerte}. Por cada millón adicional de
        publicidad, las ventas aumentan en promedio 4{,}43 millones.
      \end{alertblock}
    \end{column}
  \end{columns}
\end{frame}

\seccion{3.6 Validación de modelos}

\begin{frame}{Coeficiente de determinación $r^{2}$}
  \begin{block}{Qué significa}
    \small
    $r^{2}$ es la \destacar{proporción de la variabilidad de $y$ que el modelo
    explica}. Se expresa como porcentaje y va de 0 a 1.
  \end{block}

  \vspace{0.2cm}
  \begin{exampleblock}{En el ejercicio}
    \footnotesize
    $r^{2}=0{,}9682$: el 96{,}8\% de la variación en las ventas se explica
    por la inversión en publicidad. El 3{,}2\% restante se debe a otros
    factores no incluidos en el modelo.
  \end{exampleblock}

  \vspace{0.15cm}
  \begin{alertblock}{Un $r^{2}$ alto no valida el modelo por sí solo}
    \footnotesize
    Puede haber un $r^{2}$ alto con una relación claramente curva, o con datos
    dominados por un solo valor atípico. Siempre hay que \destacar{graficar}.
  \end{alertblock}
\end{frame}

\begin{frame}{Cómo validar el modelo}
  \small
  \begin{enumerate}
    \item \azul{Graficar la dispersión} antes de ajustar: ¿la nube es
          realmente lineal?
    \item \azul{Revisar los residuos} $e_i=y_i-\hat{y}_i$: deben repartirse al
          azar alrededor de cero, sin patrones ni curvatura.
    \item \azul{Detectar atípicos} e influyentes: un punto lejano puede
          arrastrar toda la recta.
    \item \azul{Probar la significancia} de la pendiente:
          $H_0:\beta_1=0$ contra $H_1:\beta_1\neq 0$. Si no se rechaza, la
          variable $x$ no aporta.
    \item \azul{No extrapolar}: la recta solo es creíble dentro del rango de
          $x$ observado.
  \end{enumerate}

  \vspace{0.2cm}
  \begin{alertblock}{Sobre el punto 5}
    \footnotesize
    Predecir las ventas con una inversión de 60 millones a partir de datos
    entre 2 y 14 no tiene respaldo: nada garantiza que la relación siga siendo
    lineal fuera de ese rango.
  \end{alertblock}
\end{frame}

\begin{frame}{Ejercicio 2 — en clase}
  \begin{alertblock}{Individual (9 minutos)}
    Con el modelo $\hat{y}=26{,}42+4{,}43\,x$ del ejercicio anterior:

    \vspace{0.25cm}
    \begin{enumerate}
      \item Calcule el valor estimado para $x=8$ y el residuo correspondiente
            (el valor observado fue $y=68$).
      \item Interprete $b_0=26{,}42$ en el contexto del problema.
      \item ¿Sería válido predecir las ventas con $x=60$? Justifique.
      \item ¿Puede concluirse que la publicidad \emph{causa} las ventas?
    \end{enumerate}
  \end{alertblock}
\end{frame}

\begin{frame}{Ejercicio 2 — solución}
  \small
  \begin{block}{Punto 1}
    \footnotesize
    $\hat{y}(8)=26{,}42+4{,}43(8)=\azul{61{,}86}$ \quad y el residuo es
    $e=68-61{,}86=\azul{6{,}14}$: ese mes se vendió más de lo que el modelo
    predecía.
  \end{block}

  \vspace{0.1cm}
  \begin{exampleblock}{Puntos 2 y 3}
    \footnotesize
    $b_0$ sería la venta estimada sin publicidad: 26{,}42 millones. Es
    interpretable aquí porque $x=0$ es un escenario posible, aunque queda
    fuera del rango observado.

    \vspace{0.1cm}
    Con $x=60$ \destacar{no}: sería extrapolar muy lejos del rango de 2 a 14.
  \end{exampleblock}

  \vspace{0.1cm}
  \begin{alertblock}{Punto 4}
    \footnotesize
    \destacar{No.} La correlación es fuerte, pero pudo haber una tercera
    variable (temporada, crecimiento del mercado) que moviera ambas.
  \end{alertblock}
\end{frame}

\seccion{Repaso para el Examen 3}

\begin{frame}{Mapa de la Unidad 3}
  \scriptsize
  \begin{center}
  \begin{tabular}{@{}llp{6.6cm}@{}}
    \toprule
    \textbf{Tema} & \textbf{Sesión} & \textbf{Lo esencial} \\
    \midrule
    Concepto & 13 & $H_0$ lleva la igualdad; $H_1$ es lo que se prueba \\
    Tipos & 13 & Bilateral ($\neq$) o unilateral ($>$, $<$) \\
    Errores & 13 & Tipo I ($\alpha$) y tipo II ($\beta$); potencia $1-\beta$ \\
    Valor $p$ & 13 & Si $p\le\alpha$ se rechaza $H_0$ \\
    Pasos & 14 & Plantear, fijar $\alpha$, región, calcular, concluir \\
    Estadísticos & 14 & $z$ con $\sigma$; $t$ con $s$; proporción con $p_0$ \\
    Software & 14 & Leer $p$, $gl$ y el IC de la salida \\
    Correlación & 15 & $r\in[-1,1]$; no implica causalidad \\
    Regresión & 15 & $\hat y=b_0+b_1x$; $r^{2}$ explica la variabilidad \\
    \bottomrule
  \end{tabular}
  \end{center}

  \vspace{0.1cm}
  \begin{alertblock}{Examen 3 — semana 16}
    \footnotesize
    \destacar{23 al 28 de noviembre}, en Aulas Virtuales (Moodle), presencial,
    individual. Vale el \destacar{33,4\%}. Traiga calculadora, formulario y
    tablas de $Z$ y $t$.
  \end{alertblock}
\end{frame}

\begin{frame}[fragile]{Correlación y regresión, en R}
  \scriptsize
\begin{lstlisting}[style=Rmini]
x <- c(2, 4, 5, 7, 8, 10, 12, 14)          # inversion en publicidad
y <- c(35, 40, 52, 55, 68, 70, 78, 88)     # ventas

cor(x, y)                 # coeficiente de Pearson
modelo <- lm(y ~ x)       # se lee "y en funcion de x"
coef(modelo)              # b0 y b1
summary(modelo)$r.squared # coeficiente de determinacion
predict(modelo, data.frame(x = 9))         # prediccion

plot(x, y, pch = 19, main = "Ventas segun publicidad")
abline(modelo, col = "tomato", lwd = 2)    # la recta ajustada
plot(fitted(modelo), residuals(modelo))    # validacion: residuos
\end{lstlisting}
\begin{lstlisting}[style=Rout]
[1] 0.9839578
(Intercept)           x
  26.419149    4.429787
[1] 0.9681725
       1
66.28723
\end{lstlisting}

  \vspace{-0.12cm}
  \begin{alertblock}{Lo que R no le va a decir}
    \tiny
    Que $r = 0{,}98$ \destacar{no prueba} que la publicidad cause las ventas;
    que predecir en $x = 40$ sería extrapolar fuera del rango observado; y que
    con $n = 8$ cualquier conclusión es frágil. \azul{Esos tres juicios son
    suyos}, y son los que se evalúan en el Examen 3.
  \end{alertblock}
\end{frame}

\begin{frame}{Cierre del curso}
  \begin{block}{El recorrido completo}
    \footnotesize
    \begin{itemize}
      \item \azul{Unidad 1:} describir y contar lo incierto; medir con
            probabilidad; elegir una muestra.
      \item \azul{Unidad 2:} modelar el azar con distribuciones; entender que
            $\bar X$ es aleatoria; estimar con intervalos.
      \item \azul{Unidad 3:} decidir con evidencia; relacionar y predecir.
    \end{itemize}
  \end{block}

  \vspace{0.2cm}
  \begin{alertblock}{La idea que resume el curso}
    \footnotesize
    Nunca vamos a conocer la población. Lo que la estadística inferencial
    ofrece no es certeza, sino \destacar{una medida honesta de cuánta
    incertidumbre queda} después de mirar los datos. Esa medida es lo que
    permite decidir sin engañarse.
  \end{alertblock}
\end{frame}

\begin{frame}{Ejercicios propuestos}
  \vspace{-0.1cm}
  \begin{beamercolorbox}[wd=\textwidth,sep=0.10cm]{block body}
    \centering\scriptsize
    \textbf{1--12:} datos del ejercicio 1 ($\sum x=62$, $\sum y=486$,
    $\sum xy=4287$, $\sum x^{2}=598$, $\sum y^{2}=31906$, $n=8$).
  \end{beamercolorbox}
  \listatiny
  \vspace{0.15cm}
  \begin{columns}[T,onlytextwidth]
    \begin{column}{0.47\textwidth}
      \begin{enumerate}\setlength{\itemsep}{0.16cm}
        \item $\bar{x}$
        \item $\bar{y}$
        \item $b_1$
        \item $b_0$
        \item Ecuación de la recta
        \item $r$
        \item $r^{2}$ en porcentaje
        \item $\hat{y}$ para $x=9$
        \item $\hat{y}$ para $x=6$
        \item Residuo en $x=8$ si $y=68$
      \end{enumerate}
    \end{column}
    \begin{column}{0.50\textwidth}
      \begin{enumerate}\setlength{\itemsep}{0.16cm}
        \setcounter{enumi}{10}
        \item Interprete $b_1$
        \item ¿Es válido predecir con $x=60$?
        \item Rango de valores posibles de $r$
        \item $r=-0{,}85$: ¿fuerza y sentido?
        \item $r=0{,}12$: ¿fuerza?
        \item ¿$r$ alto implica causalidad?
        \item ¿Qué mide $r^{2}$?
        \item ¿Qué se revisa en los residuos?
        \item $H_0$ y $H_1$ para la significancia de la pendiente
        \item Si no se rechaza $H_0$ del punto 19, ¿qué se concluye?
      \end{enumerate}
    \end{column}
  \end{columns}
\end{frame}

\begin{frame}{Respuestas}{Autoevaluación de los ejercicios propuestos}
  \listatiny
  \vspace{0.15cm}
  \begin{columns}[T,onlytextwidth]
    \begin{column}{0.47\textwidth}
      \begin{enumerate}\setlength{\itemsep}{0.16cm}
        \item $7{,}75$
        \item $60{,}75$
        \item $4{,}4298$
        \item $26{,}4191$
        \item $\hat y=26{,}42+4{,}43x$
        \item $0{,}9840$
        \item $96{,}82\%$
        \item $66{,}29$
        \item $53{,}00$
        \item $6{,}14$
      \end{enumerate}
    \end{column}
    \begin{column}{0.50\textwidth}
      \begin{enumerate}\setlength{\itemsep}{0.16cm}
        \setcounter{enumi}{10}
        \item Por cada millón más de publicidad, las ventas suben 4{,}43
        \item No: extrapolación fuera del rango 2--14
        \item $-1\le r\le 1$
        \item Fuerte y negativa
        \item Muy débil
        \item No
        \item La proporción de variabilidad de $y$ explicada
        \item Que se repartan al azar, sin patrón
        \item $H_0:\beta_1=0$ vs. $H_1:\beta_1\neq 0$
        \item Que $x$ no aporta a explicar $y$
      \end{enumerate}
    \end{column}
  \end{columns}
\end{frame}

\begin{frame}{Referencias}
  \footnotesize
  \begin{itemize}
    \item[{[1]}] Díaz Rodríguez, M. (2022). \textit{Estadística inferencial
      aplicada}. Universidad del Norte.
      \texttt{https://elibro.net/es/lc/uniminuto/titulos/221614}
    \item[{[2]}] Llinás Solano, H. (2014). \textit{Introducción a la
      estadística con aplicaciones en ciencias sociales}. Universidad del
      Norte. \texttt{https://elibro.net/es/lc/uniminuto/titulos/70062}
    \item[{[3]}] Contento Rubio, M. R. (2019). \textit{Estadística con
      aplicaciones en R}. Editorial Utadeo.
      \texttt{https://elibro.net/es/lc/uniminuto/titulos/220926}
  \end{itemize}

  \vspace{0.3cm}
  \begin{center}
    {\Large\color{uniAzul}\textbf{Éxitos en el Examen 3}}
  \end{center}
\end{frame}

\end{document}
